\documentclass[a4paper,11pt]{kth-mag}
\let\ifpdf\relax
\usepackage[T1]{fontenc}
\usepackage{textcomp}
\usepackage{lmodern}
%\usepackage[latin1]{inputenc}
\usepackage[swedish,english]{babel}
\usepackage{modifications}
\usepackage{amsmath}
\usepackage{ifpdf}
\usepackage[breaklinks=true]{hyperref}
\usepackage{breakcites}
\usepackage{pdfpages}
\newcommand{\textunderscript}[1]{$_{\text{#1}}$}
\usepackage[utf8]{inputenc}



%\usepackage{appendix} 
\title{Rad Cube}
\subtitle{Control of reaction wheel balanced inverted pendulum}
\foreigntitle{Rapporttitel på svenska}
\author{Mikael Sjöstedt \\ Alexander Ramm}
\date{May 2015}
\blurb{ Bacherlor's Thesis in Mechatronics \vspace{1em} \\
\begin{tabular}{ll} 
Supervisor: Daniel Frede (Is this correct?)&  \\
Examiner:Martin Edin Grimheden \\ 
Approved: & TBA 2015-month-day
\end{tabular} }
\trita{TRITA xxx yyyy-nn}


\begin{document}

%\includepdf[pages={1}]{kth-cover.pdf}
\clearpage

\frontmatter
\pagestyle{plain}
%\removepagenumbers
\pagenumbering{roman}
\maketitle
\selectlanguage{english}
\begin{abstract}
\addcontentsline{toc}{chapter}{Abstract}
Today the use of automated control is growing in a rapid pace and is being implemented more and more in consumer related products. ?This thesis is about implementing automated control and balance a simple construction using a reaction wheel commonly used in satellites? 
\\ To be filled in:
\\Problem ( 
\\Approach (On the premiss of constructing a small box shaped robot "Items" with desired properties were chosen....)
\\Results ( Hopefully something that looks cool ) 
\\ Conclusion
\\ \\
(The thesis is written in English and it has both a Swedish and an English \textit{Abstract} page. An English thesis has English \textit{Cover} and \textit{Title} pages. This template, which has no bookmarks or other automation features, defines the layout of a bachelor thesis in Machine Design (MF123X) or Industrial Engineering (MF105X).  The chapter structure that is presented here is just an example.
The abstract is maximum one page long. The abstract page is followed by a blank page. The preceding title page, which preferably contains an illustrative picture, is also followed by a blank page. All chapters start at the top of a right hand page, i.e. a page with an odd number.)
\end{abstract}
\cleardoublepage
\begin{foreignabstract}{swedish}
\addcontentsline{toc}{chapter}{Sammanfattning}
Som abstract fast på engelska
\\
Examensarbetet skrivs på engelska och har alltid både en svensk och en engelsk sammanfattningssida. Denna skrivmall, som inte har ``bookmarks'' eller några avancerade ``features'', definierar layouten för ett examensarbete i Maskinteknik (MF123X) eller Industriell ekonomi (MF105X).
I detta kapitel sammanfattas examensarbetet. Sammanfattningens omfattning är högst en sida. Sammanfattningen åtföljs av en blank sida. Den föregående titelsidan, som lämpligen också innehåller en illustrativ bild, följs också av en blank sida. Alla kapitel börjar längst upp på en högersida, dvs på ett udda sidonummer.
\end{foreignabstract}
\clearpage
\chapter*{Preface}
\addcontentsline{toc}{chapter}{Preface}
Detta är rätt ställe att tacka för hjälp, råd, samarbete och inspiration för det presenterade projektet. Detta kapitel är valfritt. 
Förordet avslutas lämpligtvis med de båda raderna Namn och Plats, månad och år. \\ \\
% Here goes credits to Daniel Frede, Staffan, Assarna som tog sig tid, eventuella maskiner som inte strulat
\begin{flushright}Alexander Ramm \\Mikael Sjöstedt \\ KTH, månad, 2015 \end{flushright}


%\clearpage


\cleardoublepage
\addcontentsline{toc}{chapter}{Contents}
\printindex
\tableofcontents*

\cleardoublepage
\chapter*{Nomenclature}
\addcontentsline{toc}{chapter}{Nomenclature}

\section*{Symbols}
\noindent{}\begin{tabular}{@{}p{2.5cm}l}
\textbf{Symbol} 	& \textbf{Description} \vspace{.5em} \\
$E$ 		& Elasticity module (Pa) \\
$r$		& Radius (m) \\
$t$		& Thickness (m) \\
L			& Lagrange ( fixa) \\
$\theta$		& Kubvinkel \\
$\phi$		& Flywheel angle \\
Q och q		& Lagrange operator \\
$E_k	$		& Kinetic energy \\
$E_p$		& Potential Energy \\
$I_c$		& Inertia of the cube\\
$I_f$		& Inertia of the flywheel\\
$M_c$		& Total mass of the cube \\
$i$			& Current\\
$K_t$		& Torque constant\\
%$_{E\{emf}}$	& Induced voltage? \\

E\textunderscript{emf} 	& Induced voltage \\
K\textunderscript{emf} 	& Induced voltage constant ??? \\

U			& Voltage over motor ?\\
$R_m	$		& Motor resitance \\
$\eta_m$		& Motor efficiency\\	
$\ldots$	& $\ldots$ 
\end{tabular}

\section*{Abbreviations}
\noindent{}\begin{tabular}{@{}p{2.5cm}l}
\textbf{Abbreviation} 	& \textbf{Description} \vspace{.5em} \\
CAD			& Computer Aided Design \\
CAE			& Computer Aided Engineering\\
PLM			& Product Lifecycle Management\\
PWM			& Pulse With Modulation\\
DOF			& Degrees of freedom\\

$\ldots$		& $\ldots$ 
\end{tabular}
\cleardoublepage

\mainmatter
\pagestyle{newchap}

\chapter{Introduction}
Beskriv bakgrund, syfte, avgränsning för det utförda examensarbetet, samt även de  metod(-er)som använts för att lösa uppgiften. Skriv exempelvis: Detta kapitel beskriver bakgrund, syfte, avgränsning och metod för det utförda examensarbetet.

This chapter decribes the backgroud,pupose, scope, and method of the theisis. 
\section{Background}
Balancing a inverted pendulum can be a challenging task. It requires knoledge of closed loop control systems and their stability, electronics and... 

\section{Purpose}

(instert RQ here)"What improvements can be made to the control system to better cope with an applied external force to the cube"  \\

\noindent{}The purpose was to examine what affected the stability of the mechanical system, (speciellt) the control system. Many parameters goes in to the state space control system and their affect on the mechanical system behaviour is not trivial. The purpose was to detirmine if some of the parameters had an extra importance and if the conclutions were applicable on other mechanical systems. Sample frequency and clock frequency were other parametesr which where to be examined aswell. Also the results were to be contributed to the open source comunity. All results is availible online, open source (referens till MIT licens här), on GitHub (ref till github länk)(summer 2015).\\  


Figurer i avhandlingen skall centreras och figurtexten skrivs centrerat. Figurerna numereras Fig. \ref{fig:kth_logo}, Fig. 1.2, och så vidare. 
\begin{figure}[htbp]
\begin{center}
\includegraphics{KTH_ITM_logo}
\caption{Logo of the ITM school.}
\label{fig:kth_logo}
\end{center}
\end{figure}

\noindent{}Tabellrubriker skrivs centrerat. Tabeller numreras Table \ref{tab:defaulttable}, Table 1.2, och så vidare. Tabelltext skrivs centrerat. Nedan visas ett exempel på en tabell med tabellrubrik.
\begin{table}[htdp]
\caption{Table header}
\begin{center}
\begin{tabular}{cc}
\toprule
Rubrik & Rubrik \\
\midrule
råd 1, kolumn 1 & råd 1, kolumn 2 \\
råd 2, kolumn 1 & råd 2, kolumn 2 \\
$\ldots$	       & $\ldots$ \\
råd n, kolumn 1 & råd n, kolumn 2 \\
\bottomrule
\end{tabular}
\end{center}
\label{tab:defaulttable}
\end{table}%


\begin{itemize}
\item Onumrerade listor skrivs lämpligtvis med i punktform (``bullets'')
\end{itemize}

\begin{enumerate}
\item Numrerade listor numreras 1,2,3 etc.
\end{enumerate}

\noindent{}Ekvationer centras. Varje ekvation numreras konsekutivt genom hela rapporten med arabiska siffror inom parantes: eq. \eqref{eq:testequation}, (1.2), etc. Ekvationsnumret högerjusteras på raden.
\begin{equation}
R= \frac{N_F}{N_F-N_R} \label{eq:testequation}
\end{equation}

\noindent{}Referenser bör göras med namn och år, som exempel, \cite{Angell1999} hävdar att $\ldots$, eller ickelinjär finita element analys förutsätter såväl specialkunskaper som erfarenhet \cite{DeBorst2012}. Observera också lilla skillnaderna mellan referenser till tidskrifter \cite{VanWezel2001} och \cite{Angell1999}, en bok \cite{DeBorst2012} och en konferens \cite{Eppinger2001}. Ett exempel på en Internetreferens är \cite{StandishGroup1994}. Huvudrubriken för referenskapitlet behandlas som en helt vanlig huvudrubrik. \\


\noindent{}\emph{Angående forskningsfrågor och vetenskaplig frågeställning}\\
Beskriv, med utgångspunkt i den bakgrund som presenterats i föregående avsnitt, den frågeställning som rapporten undersöker, diskuterar och drar slutsatser runt. Själva frågeställningen ska beskrivas i 1-2 meningar som tydligt pekar ut vad som ska mätas inom de begränsningar och med de metoder som presenteras i efterföljande avsnitt. Med andra ord, undvik vaga formuleringar (``Går det att$\ldots$''), kriterier som egentligen är funktioner (``Ska kunna$\ldots$'') eller subjektiva formuleringar (``$\ldots$ska vara bra/acceptabel/tillräckligt snabb/etc.'')
\section{Scope}
Det är god sed att definiera och beskriva projektets/uppgiftens  avgränsning i det introducerande kapitlet.

The scope were to examine the parameters, of a state space controller and other electronical desitions, affect on the behaviour of a balancing "1-DOF" inverted pendulum.
\chapter{Method}
Den metod eller de metoder, som huvudsakligen används för att angripa den uppgift eller det problem som definieras ovan, kan antingen definieras i introduktionskapitlet eller förklaras mera ingående i ett följande metodkapitel.\\

The engineering tast was devided into smaller problems (In med punktlista?). (Behandlede) seperatly to later asseble the entire system. The choosen subproblems where: 
\begin{itemize}
\item Construction
\item Motor Control
\item Sensor Reading
\item System Control
\end{itemize}

\section{Contruction}
The main contruction problem where deciding the size of the cube and reaction wheel. A too big reaction wheel for the motor has a large affect on the cubes abillity to balance. The problem where (uppställt) with Newtionial mechanics.
 
\section{Motor and Motor Control}
The motors nominal and stall torque are very important for the system blaha
  
\section{Sensor Reading}
The IMU's parameters and filtering of the signals

\section{System Control}
The choosen conrol methot where state space. The problem in to linareise and discretise with good enough precition.


\chapter{Theory}
\emph{Den teoretiska fördjupningen är en sammanfattning av tillgänglig kunskap och resultat från forskning som tidigare har utförts inom examensarbetets område. Detta kapitel presenterar den teoretiska referensramen som utgör utgångspunkten för den utförda forskningen, produktutvecklingen eller konstruktionsuppgiften.}
\\ Här måste det köttas på mycket inför tisdag 9/3 \\

Balancing 1-DOF inverted penelum type strucktures using reaction wheels is no new concept, and became common with the introduction to cheap microcontrollers. A lot work has been done on the topic but it's still no easy task due tu the instabillity of pendulums. The latest development is on "2-DOF" pendelum structures unsing multiple orthogonal reaction wheels. This method is commonly used to rotate sattelites and maintaining their attitude so incerease performance and allign solar panels. 
Allso creating transversal movement using only the reaction wheels is a recent topic for research ie not only changeing direction of something but actually moving it. This is of course impossible in orbit, but could possibly be usefull for land/sea based machines to overcome various obstacles without a seperate system for balanceing and movement. 
\chapter{Demonstrator}
\emph{Detta kapitel beskriver både den utvecklade demonstratorn och den aktuella arbetsprocessen som demonstartorn utvecklats enligt, dvs resultatet och vägen dit.}


\section{Problem Formulation}
Beskriv din problemställning för demonstratorn.\\
Create ze cube, balance it, go home
The engineering problem where to build a cube that, using a reaction wheel, could balance on it-s edge.
\section{State space model}
To create a state-space model the physical model has to be translated to a mathematical model. To do this, \emph{Euler-Lagrange} equations is used where a system in motion can be described by:
\begin{equation}
Q_i=\frac{d}{dt}\left(\frac{\partial L}{\partial \dot{q_i}}\right)-\left(\frac{\partial L}{\partial q_i}\right)
\end{equation}

In this case the cube's angular momentum (Rätt ordval?) is counteracted by the flywheel and the system can be then be written as follows

%\begin{equation}
\begin{equation} \label{eq:positiveL}
M_a=\frac{d}{dt}\left(\frac{\partial L}{\partial \dot{\theta}}\right)-\left(\frac{\partial L}{\partial \theta}\right)
\end{equation}
\begin{equation} \label{eq:negativeL}
-M_a=\frac{d}{dt}\left(\frac{\partial L}{\partial \dot{\phi}}\right)-\left(\frac{\partial L}{\partial \phi}\right)
\end{equation}



Whereas $\theta$ represents the angle of the cube and $\phi$ is the position of the flywheel. \\
The Lagrange equation is derived from the difference in kinetic energy and potential energy of the cube
\begin{equation} \label{eq:Lagrange}
L = E_k - E_p
\end{equation}

\begin{align}
E_k = \frac{I_c \cdot \dot{\theta^2}}{2} + \frac{I_f \cdot \dot{\phi^2} }{2} \\
E_p = \frac{M_c \cdot g \cdot l \cdot \cos \theta}{\sqrt{2}} 
\end{align}
Equation \eqref{eq:negativeL} and \eqref{eq:positiveL} with \eqref{eq:Lagrange} gives
\begin{equation} \label{eq:negativeL2}
I_k \cdot \ddot{\theta} - \frac{M_c \cdot g \cdot l \cdot \sin \theta }{\sqrt{2}}  = -M_a
\end{equation}
\begin{equation} \label{eq:postiveL2}
I_s \cdot \ddot{\phi} = M_a
\end{equation}
It is now plain to see that (Jävligt slentriant skrivet) $M_a$ is the torque (utövat ??) needed by the flywheel which is wielded by the motor torque $\tau$, it can be described by a relation between a motor constant and the current flowing through the motor.
\begin{equation}
\tau = K_t \cdot i_m
\end{equation}
The current can be described by the voltage (över motorn?)
\begin{equation}
\tau = K_t \cdot \frac{U-E_{\text{emf}} }{R_m}
\end{equation}
The induced voltage is related to the 'random constant´and the rotor rotation
\begin{align}
E_{\text{emf}} = K_{\text{emf}} \cdot \dot{\phi_r} \\
\phi_r = \dot{\phi} - \dot{\theta} 
\end{align}
\begin{equation} \label{eq:tau}
\tau = \frac{K_t}{R_m} U - \frac{K_t K_{\text{emf}} }{R_m} \dot{\phi} + \frac{K_t K_{\text{emf}} }{R_m} \dot{\theta}
\end{equation}
The torque on the (hur fan blir det här bra ? ) flywheel can then be described with the efficiency of the motor
\begin{equation} \label{eq:eff}
M_a = \tau \cdot \eta_m
\end{equation}
Based on equation \eqref{eq:negativeL}, \eqref{eq:positiveL} and \eqref{eq:eff} the system can be described by
\begin{align} 
\ddot{\theta} = -\frac{K_t \eta_m}{R_m I_c} U + \frac{K_t K_{\text{emf}} \eta_m}{R_m I_c} \dot{\phi} - \frac{K_t K_{\text{emf}} \eta_m}{R_m I_c} \dot{\theta} + \frac{Mt g l }{\sqrt{2} I_c} \sin \theta \label{thetadotdot}
\\
\ddot{\phi} = \frac{K_t \eta_m}{R_m I_f} U + \frac{K_t K_{\text{emf}} \eta_m}{R_m I_f} \dot{\phi} - \frac{K_t K_{\text{emf}} \eta_m}{R_m I_f} \dot{\theta} 
\label{phidotdot}
\end{align} 

based on equations \eqref{thetadotdot} and \eqref{phidotdot} the system can be described with a state space model with a states $x^T = [\theta, \dot{\theta}, \dot{\phi}$]. The system is hence described by
\begin{equation}
\dot{x} = Ax + Bu
\end{equation} 
where \\
\begin{center}
$A =\begin{bmatrix}
0 & 1 & 0 \\
\frac{Mt g l }{\sqrt{2} I_c} & - \frac{K_t K_{\text{emf}} \eta_m}{R_m I_c} & \frac{K_t K_{\text{emf}} \eta_m}{R_m I_c} \\ 
0 & \frac{K_t K_{\text{emf}} \eta_m}{R_m I_f} & -\frac{K_t K_{\text{emf}} \eta_m}{R_m I_f}
\end{bmatrix}$

$B = \begin{bmatrix}
0 \\ 
-\frac{K_t \eta_m}{R_m I_c} \\
\frac{K_t \eta_m}{R_m I_f}
\end{bmatrix} $
\end{center}


\section{Software}
Beskriv hur din mjukvara fungerar. Använd bl.a. flödesscheman för att åskådliggöra programmets struktur.

\section{Electronics}
Beskriv din elektroniska konstruktion. Använd figurer och förenklade blockschema. Motivera dina lösningar.
\\ Sensors
\\ Motor
\\ arduino
\\ Motor control


\section{Hardware}
Beskriv din mekaniska konstruktion (om du har någon)
The motor is fixed through the middle wall in the cube, the shaft on one side and the body on the other. The flywheel is dicrectly mounted to the motor shaft. All other componets are mounted on the motor-body side of the cube.
\\ Basic construction

\section{Results}
Beskriv resultatet.


\chapter{Discussion and conclusions}
\emph{I detta kapitel diskuteras och sammanfattas de resultat som presenterats i föregående kapitel. Sammanfattningen baseras på en resultatanalys och syftar till att svara på den fråga eller de frågor som formuleras i kapitel i.}

\section{Discussion}
Bla bla bla

\section{Conclusions}
Bla bla bla


\chapter{Recommendations and future work}
\emph{I detta kapitel ges rekommendationer for mera detaljerade lösningar och/eller framtida arbete.}

\section{Recommendations}
Lorem ipsum dolor sit amet, consetetur sadipscing elitr, sed diam nonumy eirmod tempor invidunt ut labore et dolore magna aliquyam erat, sed diam voluptua. At vero eos et accusam et justo duo dolores et ea rebum. Stet clita kasd gubergren, no sea takimata sanctus est Lorem ipsum dolor sit amet.

\section{Future work}
An extention of the project would be balancing the cube not only in it's edge but it's corner. To achieve this multiple reaction wheels must be used and a more complicated control system due to changes in moment of inertia caused by angular velocities in the other reaction wheels.

%\bibliography
\cleardoublepage
\bibliography{FiM_references}
\bibliographystyle{apalike-url}

\cleardoublepage
\appendix
\addtocontents{toc}{\protect\contentsline {part}{Appendices}{}{}}


\chapter{Additional information} \label{appA}

\chapter{Proofs} \label{appB}

\cleardoublepage   
\cleartoverso %force back cover to be "left" page
%\includepdf[pages={2}]{kth-cover.pdf}


\end{document}

